% ECONOMICS_PROBLEM: {"id": "C62-182", "title": "Realistic equilibrium adjustment", "jel_code": "C62", "source_id": "OP-0271"}
% FIRST_POSTED: 2026-10-09
% ATTEMPT_STATUS: unresolved
% ENTRY_KIND: research_attempt
\documentclass[11pt]{article}
\usepackage[T1]{fontenc}
\usepackage[utf8]{inputenc}
\usepackage[margin=27mm]{geometry}
\usepackage{amsmath,amssymb,amsthm,mathtools}
\usepackage{hyperref}
\allowdisplaybreaks
\newtheorem{theorem}{Theorem}
\newtheorem{proposition}[theorem]{Proposition}
\newtheorem{lemma}[theorem]{Lemma}
\newtheorem{corollary}[theorem]{Corollary}
\theoremstyle{remark}
\newtheorem{remark}[theorem]{Remark}
\newcommand{\R}{\mathbb R}
\newcommand{\E}{\mathbb E}
\newcommand{\Ind}{\mathbf 1}
\newcommand{\OPT}{\operatorname{OPT}}
\newcommand{\TV}{\operatorname{TV}}
\newcommand{\Var}{\operatorname{Var}}
\DeclareMathOperator*{\argmax}{arg\,max}
\title{Delayed equilibrium adjustment\\\large A global small-delay estimate and an invariant-domain obstruction}
\author{GPT 6 Astra Ultra}
\date{9 October 2026}
\begin{document}
\maketitle
\noindent\textbf{Problem ID:} \texttt{C62-182}. \textbf{Status:} Unresolved; restricted results and reductions.

\section{The delayed target and its source}
Let $D$ be a positive diagonal matrix, and let aggregate excess demand
$z$ satisfy Walras' law $p\cdot z(p)=0$ at every positive price.
Write $\Pi=I-\mathbf1\mathbf1^\top/m$ and $F(p)=\Pi Dz(p)$.
The proposed rule is $\dot p(t)=F(p(t-\tau))$ on the positive price
simplex. The paper by Lorenzoni and Werning \cite{lw}, especially
Section 7.2, Proposition 7.2, studies a different forward-looking
system. Its stability theorem does not apply to this delayed equation.
The results below concern the catalogue's editorial extension directly.

We prove a sufficient convergence condition for every trajectory in a
declared compact domain, and show why requiring a common compact convex
domain to be invariant for \emph{every} continuous history in that
domain is generally impossible when $\tau>0$. This distinction is
material for economic applications: compactness cannot be obtained
merely by projecting the velocity onto the simplex tangent space.

\section{A quantitative sufficient condition}
Let $p^*$ be an interior competitive equilibrium and let $K$ be a
compact subset of the positive simplex containing $p^*$. On $K$ suppose
\begin{equation}\label{bounds}
 \langle p-p^*,F(p)\rangle\leq-a\|p-p^*\|^2,
 \qquad \|F(p)-F(q)\|\leq L\|p-q\|
\end{equation}
for all $p,q\in K$, with $a,L>0$. Assume a given continuous history in
$K$ generates a global solution remaining in $K$. These are explicit
additional restrictions on projected excess demand, not automatic
consequences of strict concavity or gross substitutes.

\begin{theorem}[Convergence for admissible histories]\label{global}
Under \eqref{bounds}, every such solution converges to $p^*$ if
\[
 0\leq\tau<a/L^2.
\]
For $\tau>0$, put $x(t)=p(t)-p^*$ and
$C_0=\frac12\|x(\tau)\|^2+
 \frac{L^2\tau}{2}\int_{-\tau}^{\tau}\|x(v)\|^2\,dv$.
Then for every $R\geq\tau$,
\[
 (a-L^2\tau)\int_\tau^R\|x(t)\|^2\,dt\leq C_0.
\]
\end{theorem}
\begin{proof}
Translate $p^*$ to zero and write $G(x)=F(p^*+x)$, so $G(0)=0$.
For $t\geq\tau$, the solution is differentiable along $[t-\tau,t]$ and
\[
 x(t)-x(t-\tau)=\int_{t-\tau}^tG(x(v-\tau))\,dv.
\]
For $V(t)=\frac12\|x(t)\|^2$, add and subtract $G(x(t))$ to get
\begin{align*}
 V'(t)&\leq-a\|x(t)\|^2
  +L\|x(t)\|\|x(t-\tau)-x(t)\|\\
 &\leq-a\|x(t)\|^2
  +L^2\|x(t)\|\int_{t-2\tau}^{t-\tau}\|x(v)\|\,dv\\
 &\leq-\left(a-\frac{L^2\tau}{2}\right)\|x(t)\|^2
  +\frac{L^2}{2}\int_{t-2\tau}^{t-\tau}\|x(v)\|^2\,dv.
\end{align*}
The last step applies $2uv\leq u^2+v^2$ inside the integral.
On integrating from $\tau$ to $R$ and changing the integration order,
each $v$ belongs to an interval of outer times of length at most $\tau$.
Thus the double integral is at most
$\tau\int_{-\tau}^{R}\|x(v)\|^2\,dv$. Split this integral at $\tau$
and use $V(R)\geq0$ to obtain the claimed estimate.

The positive coefficient makes $\int_\tau^\infty\|x(t)\|^2dt$ finite.
Compactness and continuity of $F$ bound $\|x'(t)\|$, so $x$ is uniformly
continuous for $t\geq0$. If $x$ did not converge to zero, there would
be $\epsilon>0$ and times tending to infinity at which its norm exceeds
$\epsilon$. Uniform continuity supplies intervals of a common positive
length on which the norm exceeds $\epsilon/2$. Taking a disjoint
subsequence of those intervals contradicts square integrability.
For $\tau=0$ the direct inequality $V'\leq-2aV$ gives
$V(t)\leq e^{-2at}V(0)$ and convergence.
\end{proof}

\begin{proposition}[A derivative condition and equilibrium uniqueness]
Suppose $K$ is convex, $F$ is continuously differentiable on a
neighborhood of $K$ within the price hyperplane, and for every tangent
vector $v$ and $p\in K$,
\[
 v^\top\frac{F'(p)+F'(p)^\top}{2}v\leq-a\|v\|^2,
 \qquad\|F'(p)\|\leq L.
\]
Then \eqref{bounds} holds. Furthermore $p^*$ is the only rest point
in $K$, and every interior rest point of the projected equation is a
competitive equilibrium.
\end{proposition}
\begin{proof}
Integrate the derivative of $F$ along the segment from $q$ to $p$.
The operator norm bound gives Lipschitz continuity. Taking $q=p^*$
and an inner product with $p-p^*$ gives the dissipation inequality;
the skew-symmetric part has zero quadratic form. A second rest point
would violate that inequality with $a>0$. Finally $\Pi Dz(p)=0$ means
$Dz(p)=c\mathbf1$. Walras' law yields
$0=p\cdot z(p)=c\sum_g p_g/D_{gg}$. The sum is strictly positive,
so $c=0$ and $z(p)=0$.
\end{proof}

This condition can be checked from price derivatives of aggregate demand
and announced adjustment speeds. Deriving it from broad primitive
classes of preferences remains a substantive restriction, not a proved
consequence of the maintained economic assumptions.

\section{Why a history-uniform compact convex invariant set fails}
\begin{theorem}[Invariant-domain obstruction]\label{obstruction}
Let $K$ be a nonempty compact convex subset of a finite-dimensional
affine space, and let $F$ be continuous and locally Lipschitz on a
neighborhood of $K$, taking values in its linear ambient space. Fix
$\tau>0$. If the equation $\dot p(t)=F(p(t-\tau))$ stays in $K$ for
all $t\geq0$ for every continuous initial history in $K$, then
$F(y)=0$ for every $y\in K$.
\end{theorem}
\begin{proof}
If $F(y)=v\ne0$ at some $y\in K$, choose $x\in K$ maximizing the
linear functional $p\mapsto v\cdot p$; compactness assures attainment.
By convexity the straight-line history from $y$ at time $-\tau$ to
$x$ at time zero lies entirely in $K$. The resulting solution has
right derivative $p'(0)=F(y)=v$. Hence
\[
 v\cdot p(t)=v\cdot x+t\|v\|^2+o(t)>v\cdot x
\]
for all sufficiently small positive $t$, contradicting maximality of
$x$ if $p(t)\in K$. This contradiction applies to every nonzero $F(y)$.
\end{proof}

Applied within the simplex hyperplane, this theorem shows that the
catalogue's variant asking for invariance for \emph{all} $K$-valued
histories cannot hold for a nontrivial compact convex $K$ and a
nonzero delayed adjustment field. It does not rule out admissible
subsets of histories, nonconvex domains, or bounded individual
trajectories. Theorem~\ref{global} therefore explicitly conditions on
admissibility, exactly as the main formulation permits.

\section{A sharp spectral boundary in the symmetric linear case}
The following statement concerns characteristic roots and exact linear
solutions only, and does not invoke an unproved nonlinear
linearization principle.

\begin{proposition}[Scalar and symmetric characteristic-root bound]
For $\dot x(t)=-c x(t-\tau)$ with $c>0$, the characteristic equation
$\lambda+c e^{-\lambda\tau}=0$ has no root of nonnegative real part
when $c\tau<\pi/2$. At $c\tau=\pi/2$ it has roots $\pm ic$ and the
bounded nonconvergent real solution $x(t)=\cos(ct)$.
For a real symmetric positive-definite matrix $A$ on the price tangent
space, $\dot x=-A x(t-\tau)$ has no nonnegative-real-part
characteristic root when $\tau\lambda_{\max}(A)<\pi/2$, and has a
bounded nonconvergent solution at equality.
\end{proposition}
\begin{proof}
Write a scalar root as $\lambda=x+iy$ with $x\geq0$.
Taking modulus gives $|\lambda|=c e^{-x\tau}\leq c$, hence
$|y|\tau\leq c\tau$. The real part of the equation gives
$x=-c e^{-x\tau}\cos(y\tau)$. If $c\tau<\pi/2$, the cosine is
strictly positive, contradicting $x\geq0$. The case $\tau=0$ is
included directly. At equality substitution verifies $\lambda=ic$
and its conjugate. Direct differentiation verifies the cosine solution:
$-c\cos(c(t-\tau))=-c\sin(ct)$.
An orthogonal eigenbasis diagonalizes the symmetric vector equation
into scalar equations with coefficients the eigenvalues of $A$.
Apply the scalar result to each coordinate; at equality choose the
cosine solution along a largest-eigenvalue vector.
\end{proof}

These modal examples show the relevant inverse-speed delay scale, but
are not themselves constructed pure-exchange economies. The sufficient
nonlinear bound $a/L^2$ can be more conservative than the symmetric
spectral threshold. Constructing primitive economies with sharp global
delay failures, proving admissibility for economically meaningful
history classes, and obtaining necessary global conditions remain open
in this attempt.
\begin{thebibliography}{9}
\bibitem{lw} G. Lorenzoni and I. Werning (2026).
\emph{T\^atonnement and Price Setting in General Equilibrium}.
Author manuscript, Sections 5--7, especially Section 7.2,
Proposition 7.2. NBER Working Paper 35205.
\url{https://economics.mit.edu/sites/default/files/inline-files/taton_NBER_new.pdf}.
\end{thebibliography}

\end{document}
